\scanpage{100}
\deleted{\unsure{provenant} de $\mathrm{H}^1(X,G)$, \unsure{prouver} \unsure{par} \unsure{un} \inserted{\unsure{petit}} \unsure{argument} \unsure{positif}}

On utilisera un argument de spécialisation.

\unsure{On en} \deleted{\unsure{partant de}} \inserted{\unsure{suivant}} Cartier, de la suite exacte
\[
0\unsure{{-}|\longrightarrow}\underline{\mathcal{O}}_X\xrightarrow{F}\underline{\mathcal{O}}_X\longrightarrow\underline{\mathcal{O}}_X/\underline{\mathcal{O}}_X^p\unsure{{-}|\longrightarrow}0
\]
\[
1\unsure{{-}|\longrightarrow}\underline{\mathcal{O}}_X^*\xrightarrow{F}\underline{\mathcal{O}}_X^*\unsure{{-}|\longrightarrow}\underline{\mathcal{O}}_X^*/\underline{\mathcal{O}}_X^{*p}\longrightarrow1
\]
\margin{($X$ intègre !)}
\deleted{\unsure{on tire}} on tire
\[
\mathrm{H}^1(X,\alpha_p)=\mathrm{H}^1(X,\underline{\mathcal{O}}_X)^{1-F}\simeq\mathrm{H}^0(X,\underline{\mathcal{O}}_X/\underline{\mathcal{O}}_X^p)\xrightarrow{df}\mathrm{H}^0(X,\underline{\widetilde{\Omega}}^1_X)
\]
\[
\mathrm{H}^1(X,\mu_p)=\mathrm{H}^1(X,\underline{\mathcal{O}}_X)^{1-F}\simeq\mathrm{H}^0(X,\underline{\mathcal{O}}_X/\underline{\mathcal{O}}_X^*)\xrightarrow{\frac{df}{f}}\mathrm{H}^0(X,\underline{\widetilde{\Omega}}^1_X)
\]
et si $X$ est \uline{normal} les deux flèches écrites sont injectives [\deleted{\unsure{leurs}}\inserted{en tous cas,} les noyaux s'identifient, soit aux \uline{diviseurs additifs} sur $X$ annulés par $F$, soit aux diviseurs multiplicatifs annulés par $F$].

Donc on est ramené au

\uline{Lemme}\quad $\mathrm{H}^0(X,\widetilde{\Omega}^1_{X/k})\longrightarrow\mathrm{H}^0(Y_{\overline{K}},\widetilde{\Omega}^1_{Y_{\overline{K}}/\overline{K}})$ est injectif.

La vérification est presque triviale\ldots

\uline{Corollaire}\quad $\underline{\mathrm{Pic}}^\tau_{X_{\overline{K}}/\overline{K}}\longrightarrow\underline{\mathrm{Pic}}^\tau_{\unsure{X|Y}_{\overline{K}}/\overline{K}}$ a un noyau qui est un \uline{$p$-groupe fini unipotent}, [le fait qu'il est complet résulte déjà de $\operatorname{prof}X\geq2$, grâce au \deleted{Bertini} Lefschetz-Grauert-Bertini], ``sans partie fixe''. Si $X$ est normal (pas néc. de $\dim\geq2$), \inserted{$\underline{\mathrm{Pic}}^{\unsure{0|\tau}}_{X/k}$ est propre sur $k$.}
\margin{purement infinitésimal}

Le ``sans partie fixe'' provient de la normalité; elle explique que l'on a une \uline{immersion} \uline{finie} $\underline{\mathrm{Pic}}^\tau_{X/k}\longrightarrow\prod_i\underline{\mathrm{Pic}}_{Y_i/k}$, les $Y_i$ des sections hyperplanes \unsure{rendues} normales. Donc pour prouver
\margin{99}
